\section{Discussion}\label{sec:disc}

Given who rents and who owns, the second mode in the income of older Danes living alone is produced by arithmetic. A single person living on the folkepension has, after tax, a disposable income that varies little from one pensioner to the next; a single person living on the folkepension who rents her home adds to it a tax-free benefit whose median is close to a third of the income of pensioners without it. Two thirds of the older population living alone rent, and half of the older population living alone depends on public transfers, so both groups are large, and their incomes are concentrated enough that the gap between them, 26{,}000 to 42{,}000 kroner, is wider than the interquartile range within either part of the refined group. The argument takes tenure as given and says nothing about why people rent or own, nor about how the income test of the benefit affects other income; the exercise of Section~\ref{sec:cf}, which removes the benefit from income, is an accounting statement about where the modes come from, not an estimate of what incomes would be without the benefit.

The Norwegian comparison shows what the arithmetic needs. Norway has a minimum pension that also sets the disposable income of those who receive it, and a means-tested housing allowance. What it lacks is a large group living on the minimum pension: because the earnings-related national pension exceeds the minimum for almost everyone with a working history, only four to seven per cent of older Norwegians living alone are transfer-reliant. A floor can shape an income distribution only if a large share of the group stands on it. The other countries of Section~\ref{sec:countries} show that the requirements are general. The Netherlands, the United Kingdom and France until 2014 have groups on a flat floor and a housing allowance that lifts part of them clearly above it, and their distributions have the Danish shape; Sweden, Germany, the United States, Ireland and France after 2014 lack one of the ingredients, and no second mode in theirs is due to the allowance. The French case is the closest thing to an experiment: as the allowance ceased to separate recipients from non-recipients, the second mode ceased to be found once the later waves were pooled in. The size of the recipient share matters as well: in Denmark a second transfer received by a fifth of those on the floor still splits the distribution in large samples, and one received by a tenth does not.

The results bear on the measurement of low income in three ways. The conventional at-risk-of-poverty line, 60 per cent of the median of equivalised disposable income (square-root scale, persons), was 140{,}974 kroner in 2016 and 151{,}375 kroner in 2019 in the LIS data, and in those years it falls in the valley between the two Danish modes. Among the refined transfer-reliant group, 91 and 92 per cent of those without the housing benefit are below it, and 9 and 10 per cent of those with it. The line has not always been there: non-recipients have been below it in every wave, 91 to 97 per cent of them, but the share of recipients below it has fallen from 43 and 51 per cent in 2004 and 2007 to a tenth, as the line moved from the upper mode into the valley. For older Danes who live on public transfers, being counted as at risk of poverty is thus almost entirely a matter of whether one receives the housing benefit, which in turn is almost entirely a matter of renting. A line in the valley is also unusually sensitive to its own level and to the benefit rates that fix the two modes. In Norway the same line lies above most of the transfer-reliant group, with and without the housing allowance.

As to housing, the lower mode is formed mainly by owners, who pay no rent, and the upper mode by renters, part of whose income is a subsidy towards rent. Disposable income as conventionally measured, without imputed rent for owners and before housing costs, therefore ranks renters on the pension above owners on the pension, although the renters' extra income is a partial reimbursement of a cost that the owners do not bear. This is an instance of a general point made in the literature on imputed rent and on incomes after housing costs \citep{FrickGrabka2003,FrickGrabkaSmeedingTsakloglou2010,CanberraHandbook2011,List2023}; the Danish case shows that the convention can create a whole mode of the distribution.

Finally, polarisation indices \citep{EstebanRay1994,DuclosEstebanRay2004} and mixture models of the income distribution \citep{PittauZelliJohnson2010} respond to exactly the kind of structure found here. In a population in which statutory transfers are a large part of income, such measures may track the parameters of benefit schedules more closely than anything that could be called social polarisation.

The results have limits. The data are cross-sections, so nothing is said about how individuals move between the groups; the transfer-reliant are defined by the LIS classification of income sources, which in Denmark counts civil-service pensions as public, a point that has to be handled explicitly; outside Denmark and Norway the LIS files are samples, and the other countries can be tested only with waves pooled on a common scale, which assumes that the benefit schedules keep their relative positions from wave to wave; the Irish result, whose location depends on the scale, shows that this assumption can fail. A statutory floor shapes the income distribution when a large share of a group lives on it, and a second transfer that some of them receive and others do not splits the floor in two.
